% !TEX root = ../Projektdokumentation.tex
\section{Fazit} 
\label{sec:Fazit}

\subsection{Soll-/Ist-Vergleich}
\label{sec:SollIstVergleich}
Bei einer rückblickenden Betrachtung des Demo-Abschlussprojektes, kann festgehalten werden, dass alle zuvor festgelegten Anforderungen gemäß dem Pflichtenheft erfüllt wurden.
Der während der Planungsphase erstellte Projektplan(Abschnitt \bhref{sec:Projektphasen}) konnte eingehalten werden. Es mussten aufgrund von Feedback zusätzliche Anforderungen 
während der Entwicklung umgesetzt werden, dies hat aber den Projektverlauf aufgrund geplanter Pufferzeiten nicht beachtenswert beeinflusst.

\subsection{Lessons Learned}
\label{sec:LessonsLearned}
Im Zuge des Projektes konnte die Autoren wertvolle Erfahrungen bezüglich der Planung und Durchführung von Projekten sammeln. Auch wurde deutlich wie wichtig eine sinnvolle
Planung ist und welche Einflüsse neue Anforderungen oder Änderungen von Anforderungen während der Entwicklung haben können. Die Skalierung der View-Objekte innerhalb des
Tkinter-Frameworks gestaltete sich zu Beginn etwas schwierig. Ebenso war die Bereitstellung einer Datenbankanbindung über verschiedene Endgeräte ein erstes Hindernis.
Abschließend lässt sich sagen, dass das Demo-Projekt als gute Übung für das eigentliche IHK-Abschlussprojekt diente. 

\subsection{Ausblick}
\label{sec:Ausblick}
Obwohl alle im Lastenheft definierten Anforderungen realisiert werden konnten, gibt es durchaus Ansätze für zukünftige Erweiterungsmöglichkeiten. Es könnte eine Mobile
Payment Integration erweitert werden. Ebenso könnte über eine \acs{KI}-gestützte Nachfrageprognose nachgedacht werden, welche einen strukturierteren Nachfüllprozess ermöglichen
würde. Die Anwendung könnte als  Multi-Festival-Plattform erweitert werden, um Funktionen nicht nur Festival-spezifisch bereitzustellen. 
